\documentclass[11pt,a4paper]{article}
\usepackage[margin=1in]{geometry}
\usepackage{amsmath,amssymb}
\usepackage{booktabs}
\usepackage{array}
\usepackage{hyperref}
\usepackage{xcolor}
\usepackage{listings}
\usepackage{parskip}
\usepackage{enumitem}
\usepackage{mdframed}
\usepackage{graphicx}

\definecolor{impacthigh}{RGB}{0,100,0}
\definecolor{impactmed}{RGB}{180,100,0}
\definecolor{impactlow}{RGB}{100,100,100}
\definecolor{codebg}{RGB}{245,245,245}
\definecolor{fixred}{RGB}{160,0,0}

\lstset{
    basicstyle=\ttfamily\small,
    backgroundcolor=\color{codebg},
    frame=single,
    breaklines=true,
    columns=fullflexible,
    keywordstyle=\color{blue},
    commentstyle=\color{gray},
}

\newcommand{\high}[1]{\textcolor{impacthigh}{\textbf{[HIGH]} #1}}
\newcommand{\med}[1]{\textcolor{impactmed}{\textbf{[MED]} #1}}
\newcommand{\low}[1]{\textcolor{impactlow}{\textbf{[LOW]} #1}}
\newcommand{\bug}[1]{\textcolor{fixred}{\textbf{[BUG]} #1}}
\newcommand{\file}[1]{\texttt{#1}}

\title{\textbf{IBI Graph Pretraining: Improvements and Fixes}\\
\large A systematic analysis of the current implementation with\\
concrete proposals for each component}
\author{ECG Foundation Model Project}
\date{\today}

\begin{document}
\maketitle

\begin{mdframed}
\textbf{How to read this document.}
Each improvement is tagged with its expected impact:
\textcolor{impacthigh}{\textbf{[HIGH]}} = likely to measurably improve downstream performance,
\textcolor{impactmed}{\textbf{[MED]}} = useful but secondary,
\textcolor{impactlow}{\textbf{[LOW]}} = optional polish.
\textcolor{fixred}{\textbf{[BUG]}} marks actual bugs in the current code that should be fixed regardless.
All file references are relative to \file{ibi\_graph\_model/pretraining/}.
\end{mdframed}

\tableofcontents
\newpage

% ─────────────────────────────────────────────────────────────────────────────
\section{Bug Fixes (Fix First)}
% ─────────────────────────────────────────────────────────────────────────────

\subsection{\bug{Graph built from wrong embeddings}}

\textbf{File:} \file{model.py}, \texttt{IBIGraphEncoder.forward()} lines 117--124
and \texttt{masked\_forward()} lines 179--204.

\textbf{Current code:}
\begin{lstlisting}[language=Python]
def forward(self, beats, valid_mask):
    e = self.feat_enc(beats)           # raw MLP output [B, N, D]
    e = e * valid_mask.unsqueeze(-1)
    x = self.temporal(e, valid_mask)   # temporal context
    A = build_graph(e, valid_mask, k=self.k)   # <-- graph from e, NOT x
    for layer in self.gnn:
        x = layer(x, A)
    return self.norm(x)
\end{lstlisting}

The graph adjacency matrix $A$ is built from the raw MLP embedding \texttt{e},
\emph{before} the Temporal Transformer has run. This means the K-NN topology
never sees temporal context --- two beats are linked by cosine similarity of
their raw 10-feature MLP projections, not by their contextually enriched representations.

\textbf{Fix: build the graph from the temporal encoder output} \texttt{x}:
\begin{lstlisting}[language=Python]
def forward(self, beats, valid_mask):
    e = self.feat_enc(beats)
    e = e * valid_mask.unsqueeze(-1)
    x = self.temporal(e, valid_mask)          # [B, N, D]
    A = build_graph(x, valid_mask, k=self.k)  # graph from temporal output
    for layer in self.gnn:
        x = layer(x, A)
    return self.norm(x)
\end{lstlisting}

The same issue exists in \texttt{masked\_forward()}: the graph is built from
\texttt{e} (line 198), but should be built from the temporal output of the
masked sequence. This means the GNN never propagates contextually enriched
information across similar beats.

\subsection{\bug{Reconstruction target uses online encoder, not EMA target}}

\textbf{File:} \file{model.py}, \texttt{masked\_forward()} lines 186--189.

\begin{lstlisting}[language=Python]
with torch.no_grad():
    target = enc.temporal(e, valid_mask)   # enc = self.online_encoder
    target = enc.norm(target) * valid_mask.unsqueeze(-1)
\end{lstlisting}

The teacher for masked reconstruction is the \emph{online} encoder's temporal
output (with stop-gradient), not the EMA target encoder.
This is weaker than it could be: the target representation drifts
with each gradient step and is not the stable, momentum-averaged version.
The EMA target encoder is only used for BYOL, not for masked reconstruction.

\textbf{Fix: run masked reconstruction teacher through the full EMA target encoder:}
\begin{lstlisting}[language=Python]
with torch.no_grad():
    target = self.target_encoder(beats, valid_mask)  # full EMA: MLP + temporal + GNN
    target = target * valid_mask.unsqueeze(-1)
\end{lstlisting}

This aligns with data2vec \cite{baevski2022data2vec}, where the reconstruction
target is the EMA-averaged full-encoder output, not a partial forward pass.

% ─────────────────────────────────────────────────────────────────────────────
\section{Masking Strategy}
% ─────────────────────────────────────────────────────────────────────────────

\subsection{\high{Replace uniform random masking with span/block masking}}

\textbf{File:} \file{dataset.py}, \texttt{IBIGraphCollator.\_\_call\_\_()}.

\textbf{Current:} 30\% of valid beats are randomly masked, each independently.
With K=6 neighbors in the graph, a masked beat almost always has unmasked
immediate neighbors that trivially supply its representation by interpolation.
The task is too easy.

\textbf{Problem example:} If beat $i$ is masked but beats $i{-}1$ and $i{+}1$
are unmasked, the GNN can reconstruct $i$ by simple averaging of its temporal
neighbors --- no global rhythm understanding needed.

\textbf{Fix: span masking.}
Mask contiguous spans of length drawn from a geometric distribution
(mean span length $\ell$, total mask budget 30\%):
\[
    \ell \sim \text{Geom}(p), \quad p = 0.2 \implies \mathbb{E}[\ell] = 5 \text{ beats}
\]

\begin{lstlisting}[language=Python]
def span_mask(valid_idx, mask_ratio=0.3, mean_span=5):
    n_valid = len(valid_idx)
    n_mask  = int(n_valid * mask_ratio)
    masked  = set()
    p = 1.0 / mean_span
    while len(masked) < n_mask:
        start = random.randint(0, n_valid - 1)
        span  = np.random.geometric(p)
        for k in range(start, min(start + span, n_valid)):
            masked.add(k)
            if len(masked) >= n_mask:
                break
    return masked
\end{lstlisting}

Span masking forces the model to infer missing rhythm segments from
\emph{non-adjacent} context --- exactly the kind of long-range reasoning
a graph-based model is uniquely suited for.

\subsection{\med{Curriculum masking schedule}}

Start with a low mask ratio (15\%) and increase linearly to 50\% over training.
Early in training the graph structure is poor (random weights),
so easy masking allows the encoder to first learn basic representations.
As the encoder matures, harder masking drives deeper structure learning.

\[
    r(t) = 0.15 + 0.35 \cdot \frac{t}{T}, \quad t \in [0, T]
\]

Pass the current epoch to the collator or sampler and update \texttt{node\_mask\_ratio} accordingly.

% ─────────────────────────────────────────────────────────────────────────────
\section{Pretraining Objectives}
% ─────────────────────────────────────────────────────────────────────────────

\subsection{\high{Add HRV feature prediction as auxiliary head}}

\textbf{File:} \file{signal\_utils.py} already contains \texttt{compute\_hrv\_features()}
which computes 14 window-level HRV features (SDNN, RMSSD, pNN50, LF power, HF power,
LF/HF ratio, Poincaré SD1/SD2, etc.) --- but \emph{this function is never used
during pretraining}.

\textbf{Proposed objective:}
Add a lightweight MLP head that takes the global attention-pooled representation
and predicts the HRV feature vector of that window.
The target labels are computed cheaply from the raw IBI at data-loading time.

\[
    \mathcal{L}_{\text{hrv}} = \frac{1}{|\mathcal{F}|} \sum_{f \in \mathcal{F}}
    \mathrm{smooth}\text{-}L_1\bigl(\hat{y}_f,\, y_f\bigr)
\]

where $\mathcal{F}$ are the 14 HRV features, log-transformed and standardized
across the dataset.

\textbf{Why this works:}
The encoder is forced to produce a pooled representation that is
\emph{sufficient for computing all standard HRV metrics} --- which is exactly
the representation most useful for clinical downstream tasks.
This is a form of \emph{self-prediction} grounded in decades of cardiology,
not a learned codebook.

\textbf{Combined loss:}
\[
    \mathcal{L}_{\text{total}} = \mathcal{L}_{\text{mask}} + 0.5\,\mathcal{L}_{\text{BYOL}}
    + \lambda_{\text{hrv}}\,\mathcal{L}_{\text{hrv}}, \quad \lambda_{\text{hrv}} = 0.1
\]

\textbf{Implementation sketch:}
\begin{lstlisting}[language=Python]
# In model.py
self.hrv_head = nn.Sequential(
    nn.Linear(cfg.d_model, 64), nn.GELU(), nn.Linear(64, 14)
)

# In dataset.py __getitem__
hrv_feats = compute_hrv_features(ibi_clean)
hrv_vec   = np.array([hrv_feats[k] for k in HRV_KEYS], dtype=np.float32)
# log-transform power features, standardize with precomputed stats
\end{lstlisting}

\subsection{\high{Add future beat prediction (autoregressive auxiliary)}}\label{sec:future}

Given the first $T/2$ beats of a window, predict the representation of the
next $T/2$ beats. This is a natural SSL objective for a sequential cardiac signal
and requires learning rhythm regularity and temporal dynamics.

\textbf{Implementation:}
\begin{enumerate}[leftmargin=*]
    \item Split each window into a causal half [0, $T/2$] and a target half [$T/2$, $T$].
    \item Encode the causal half with the online encoder.
    \item Encode the full window with the EMA target encoder.
    \item Add a \emph{causal prediction head} (small MLP) on top of the pooled causal representation.
    \item Loss: cosine distance between the predicted and target pooled representations.
\end{enumerate}

\[
    \mathcal{L}_{\text{pred}} = 1 - \cos\!\bigl(f_\theta(\mathbf{z}_{1:T/2}),\;
    \mathbf{z}^{\text{EMA}}_{1:T}\bigr)
\]

This requires no new data and reuses the existing EMA target encoder.

\subsection{\med{Replace BYOL with data2vec-style multi-layer target}}

\textbf{Current BYOL} targets a single pooled vector from the EMA encoder.
\textbf{data2vec} \cite{baevski2022data2vec} instead targets the
\emph{average of the top-$K$ transformer layers} from the EMA encoder,
computed on the unmasked input.

\[
    \hat{\mathbf{y}}_i = \frac{1}{K} \sum_{l=L-K}^{L} \mathbf{h}_i^{(l), \text{EMA}}
\]

This produces richer, more contextualized targets than a single final layer,
and the masked reconstruction loss is computed \emph{per-node}
(not just on masked positions from a pooled vector).
This would unify the masked reconstruction and BYOL objectives into a single
cleaner framework.

\subsection{\med{Add a within-batch contrastive term}}

BYOL has no negative examples --- it relies entirely on the EMA to prevent collapse.
Adding a small InfoNCE term across the batch provides an explicit repulsion signal:

\[
    \mathcal{L}_{\text{NCE}} = -\frac{1}{B}\sum_{i=1}^{B}
    \log \frac{\exp(\mathbf{z}_i \cdot \mathbf{z}_i^+ / \tau)}
    {\sum_{j=1}^{B} \exp(\mathbf{z}_i \cdot \mathbf{z}_j^+ / \tau)},
    \quad \tau = 0.07
\]

where $\mathbf{z}_i^+$ is the EMA target representation of the same subject.
Add with a small weight ($\lambda = 0.1$) so it does not dominate.

% ─────────────────────────────────────────────────────────────────────────────
\section{Augmentation Strategy}
% ─────────────────────────────────────────────────────────────────────────────

\subsection{\high{Add block dropout (realistic sensor loss)}}

\textbf{Current} \texttt{ibi\_dropout} removes \emph{randomly selected individual beats}
with probability 0.1.
In practice, ECG signal loss is \emph{bursty} --- you lose a contiguous segment.

\textbf{Fix: block dropout.}
Drop a contiguous segment of beats:
\begin{lstlisting}[language=Python]
def ibi_block_dropout(ibi, min_gap_beats=3, max_gap_beats=10):
    if len(ibi) <= max_gap_beats + 4:
        return ibi
    gap   = np.random.randint(min_gap_beats, max_gap_beats + 1)
    start = np.random.randint(0, len(ibi) - gap)
    return np.concatenate([ibi[:start], ibi[start + gap:]])
\end{lstlisting}

This simulates realistic sensor dropout and forces the model to learn
rhythm continuity across gaps --- directly useful for wearable ECG applications.

\subsection{\med{Add HR-trend perturbation}}

\textbf{Current \texttt{ibi\_scale}} multiplies all IBIs by a uniform constant,
shifting the mean HR but preserving the relative HRV pattern.
A more challenging and realistic augmentation is to perturb the
\emph{trend} (HR acceleration / deceleration) by adding a linear ramp:

\begin{lstlisting}[language=Python]
def ibi_trend_augment(ibi, max_slope_ms=2.0):
    """Add a linear ramp to simulate HR drift (e.g., exercise onset/offset)."""
    N     = len(ibi)
    slope = np.random.uniform(-max_slope_ms, max_slope_ms)  # ms change per beat
    ramp  = np.linspace(0, slope * N / 1000.0, N).astype(np.float32)
    return np.clip(ibi + ramp, 0.1, 3.0)
\end{lstlisting}

\subsection{\med{Add time reversal as a hard negative}}

A reversed IBI sequence has the same beat distribution but inverted dynamics
(decelerations become accelerations).
The model should \emph{not} map a forward and reversed sequence to the same representation ---
they reflect different physiological states.
Use time reversal as a \emph{hard negative} in a within-batch contrastive term
rather than as a view for BYOL.

\subsection{\low{Frequency-domain perturbation (LF/HF ratio shift)}}

Simulate sympathetic/parasympathetic balance shifts by:
(1) interpolating the IBI sequence to uniform time,
(2) applying a filter that attenuates LF ($0.04$--$0.15$ Hz) or HF ($0.15$--$0.4$ Hz) power,
(3) converting back to IBI.
This is computationally expensive but creates physiologically meaningful views
that differ in autonomic tone rather than just noise level.

% ─────────────────────────────────────────────────────────────────────────────
\section{Graph Construction}
% ─────────────────────────────────────────────────────────────────────────────

\subsection{\high{Vectorize \texttt{build\_graph()} --- remove Python loops}}

\textbf{File:} \file{model.py}, \texttt{build\_graph()}, lines 15--36.

The current implementation has a Python loop over the batch dimension and
inner loops over K-NN indices. For a batch size of 128 with 75 nodes each,
this is called twice per forward pass.

\textbf{Vectorized replacement:}
\begin{lstlisting}[language=Python]
def build_graph(x, valid_mask, k=6):
    B, N, D = x.shape
    x_n = F.normalize(x, dim=-1)                      # [B, N, D]
    sim = torch.bmm(x_n, x_n.transpose(1, 2))         # [B, N, N]

    # Mask padding nodes out of similarity
    pad = ~valid_mask                                  # [B, N]
    sim = sim.masked_fill(pad.unsqueeze(1), -1e9)
    sim = sim.masked_fill(pad.unsqueeze(2), -1e9)

    # K-NN adjacency (top-k, exclude self)
    sim_no_self = sim - torch.eye(N, device=x.device).unsqueeze(0) * 1e9
    kk  = min(k, N - 1)
    _, idx = torch.topk(sim_no_self, k=kk, dim=-1)    # [B, N, K]
    A_knn = torch.zeros(B, N, N, device=x.device)
    A_knn.scatter_(2, idx, 1.0)
    A_knn = ((A_knn + A_knn.transpose(1, 2)) > 0).float()  # undirected

    # Temporal edges (i <-> i+1 for valid nodes)
    A_temp = torch.zeros(B, N, N, device=x.device)
    A_temp[:, :-1, 1:] += torch.diag(torch.ones(N - 1, device=x.device)).unsqueeze(0)
    A_temp = A_temp + A_temp.transpose(1, 2)
    # zero out edges involving padding
    A_temp = A_temp * valid_mask.unsqueeze(1) * valid_mask.unsqueeze(2)

    # Self-loops for valid nodes
    A_self = torch.diag_embed(valid_mask.float())

    return (A_knn + A_temp + A_self).clamp(0, 1)
\end{lstlisting}

\subsection{\med{Adaptive K based on sequence length}}

With a fixed $K=6$ neighbors, short sequences (e.g., 15 beats after augmentation dropout)
are nearly fully connected, while long sequences may be too sparse.
Use $K = \max(2, \min(6, \lfloor N/8 \rfloor))$ to keep connectivity density roughly constant.

\subsection{\med{Add multi-hop edges for global rhythm context}}

The current GNN with 3 layers and $K=6$ neighbors has an effective receptive field
of at most 18 beats (3 hops $\times$ 6 neighbors) --- for 75 beats, this covers
only 24\% of the sequence from any node.

\textbf{Fix: add a global token (CLS node)} connected to all valid nodes.
This is analogous to the CLS token in BERT, but implemented as a special graph node:
\begin{itemize}[leftmargin=*]
    \item Append one extra node per batch element with a learnable embedding.
    \item Connect it to all valid nodes bidirectionally.
    \item After GNN propagation, the CLS node embedding serves as the global representation
          (replaces attention pooling).
\end{itemize}
This gives every node a 1-hop path to every other node via the CLS node,
solving the limited-receptive-field problem without increasing GNN depth.

% ─────────────────────────────────────────────────────────────────────────────
\section{Architecture Improvements}
% ─────────────────────────────────────────────────────────────────────────────

\subsection{\med{Add sinusoidal positional encoding}}

\textbf{File:} \file{model.py}, \texttt{TemporalEncoder}.

\textbf{Current:} Beat position is encoded as a raw scalar feature (feature index 9:
\texttt{position} $\in [0,1]$), passed through the MLP with all other features.
The Transformer has no dedicated positional encoding.

\textbf{Fix: add sinusoidal positional encoding} (or learnable 1-D PE) applied
before the Temporal Transformer:
\[
    \mathbf{PE}(i, 2k)   = \sin\!\left(\frac{i}{10000^{2k/d}}\right), \quad
    \mathbf{PE}(i, 2k+1) = \cos\!\left(\frac{i}{10000^{2k/d}}\right)
\]

This gives the Transformer a dedicated, rich positional signal rather than
a single scalar squeezed through a 10$\to$128 bottleneck.

\subsection{\med{Separate node-level head for masked reconstruction}}

\textbf{Current} decoder is a single linear layer $\mathbb{R}^D \to \mathbb{R}^D$.
A 2-layer MLP decoder with LayerNorm produces richer targets and is standard
in masked autoencoders (MAE):
\begin{lstlisting}[language=Python]
self.decoder = nn.Sequential(
    nn.Linear(cfg.d_model, cfg.d_model),
    nn.LayerNorm(cfg.d_model),
    nn.GELU(),
    nn.Linear(cfg.d_model, cfg.d_model),
)
\end{lstlisting}

\subsection{\low{Scale up: d\_model=256, transformer\_layers=4}}

The current \texttt{d\_model=128} with 2 Transformer layers is small for a foundation model.
With 75 nodes $\times$ 10 features, a wider model can capture richer HRV representations.
Doubling to \texttt{d\_model=256} with 4 Transformer layers and 8 attention heads
roughly quadruples parameter count ($\sim$2M $\to$ $\sim$8M) --- still lightweight
for a pretraining task.

% ─────────────────────────────────────────────────────────────────────────────
\section{Multi-Scale Pretraining}
% ─────────────────────────────────────────────────────────────────────────────

\subsection{\med{Multi-resolution windows}}

\textbf{Current:} Always 30-second windows ($\leq$75 beats at 150\,BPM).
Different HRV phenomena operate at different time scales:

\begin{itemize}[leftmargin=*]
    \item \textbf{Short (5--10\,s, $\leq$25 beats):} Respiratory sinus arrhythmia, HF power, beat-to-beat dynamics.
    \item \textbf{Medium (30\,s, $\leq$75 beats):} Standard HRV window, LF and HF bands.
    \item \textbf{Long (60--120\,s, $\leq$200 beats):} VLF power, autonomic regulation, sleep transitions.
\end{itemize}

\textbf{Simple fix:} At each training step, randomly sample a window duration
$w \sim \{10, 30, 60\}$ seconds and use a shared encoder.
This forces the encoder to be scale-aware and improves generalization
to downstream tasks with different window requirements.

% ─────────────────────────────────────────────────────────────────────────────
\section{Training Dynamics}
% ─────────────────────────────────────────────────────────────────────────────

\subsection{\med{Normalize target embeddings before MSE loss}}

\textbf{File:} \file{loss.py}, \texttt{masked\_mse\_cos()}.

\textbf{Current:} MSE is computed in raw embedding space.
If embedding magnitudes vary across dimensions, MSE is dominated by
high-variance dimensions. Normalizing both target and reconstruction
before MSE (as in data2vec) makes the loss scale-free:

\[
    \mathcal{L}_{\text{mask}} = \frac{1}{|\mathcal{M}|}
    \sum_{i \in \mathcal{M}}
    \bigl\|\bar{\mathbf{h}}_i - \bar{\mathbf{r}}_i\bigr\|^2,
    \quad
    \bar{\mathbf{v}} = \frac{\mathbf{v} - \mu}{\sigma + \epsilon}
\]

where $\mu, \sigma$ are the instance-level mean and std of the target embeddings
(per-sample normalization, not batch).

\subsection{\low{Increase EMA warmup aggressiveness}}

The current cosine EMA ramp from 0.996 $\to$ 0.9999 over 100 epochs is slow.
At epoch 1, momentum 0.996 means the target updates at $1 - 0.996 = 0.4\%$
per step --- the target and online networks are nearly identical, giving
trivial targets for masked reconstruction.

Consider starting with a lower momentum (0.9) for the first 5 epochs,
then ramping to 0.996--0.9999:
\[
    m(t) = \begin{cases}
    0.90 + 0.096 \cdot (t / t_{\text{warmup}}) & t < t_{\text{warmup}} \\
    0.996 + 0.0039 \cdot \cos\!\left(\frac{t - t_{\text{warmup}}}{T - t_{\text{warmup}}} \cdot \pi\right)^{-1} & t \geq t_{\text{warmup}}
    \end{cases}
\]

This allows the target network to diverge from the online network quickly
at the start, providing more meaningful (non-trivial) reconstruction targets.

% ─────────────────────────────────────────────────────────────────────────────
\section{Summary and Priority Order}
% ─────────────────────────────────────────────────────────────────────────────

\begin{table}[h!]
\centering
\renewcommand{\arraystretch}{1.4}
\begin{tabular}{clcc}
\toprule
\textbf{Priority} & \textbf{Improvement} & \textbf{File(s)} & \textbf{Effort} \\
\midrule
\textcolor{fixred}{\textbf{1}} & Build graph from temporal encoder output (bug) & \file{model.py} & 1 line \\
\textcolor{fixred}{\textbf{2}} & Use EMA target encoder for recon. teacher (bug) & \file{model.py} & 3 lines \\
\textcolor{impacthigh}{\textbf{3}} & Vectorize \texttt{build\_graph()} & \file{model.py} & 30 lines \\
\textcolor{impacthigh}{\textbf{4}} & Span/block masking instead of uniform & \file{dataset.py} & 20 lines \\
\textcolor{impacthigh}{\textbf{5}} & HRV feature prediction auxiliary head & \file{model.py}, \file{loss.py}, \file{dataset.py} & 1--2 hours \\
\textcolor{impacthigh}{\textbf{6}} & Future beat prediction auxiliary & \file{model.py}, \file{loss.py} & 1--2 hours \\
\textcolor{impactmed}{\textbf{7}} & Block dropout augmentation & \file{dataset.py} & 10 lines \\
\textcolor{impactmed}{\textbf{8}} & Sinusoidal positional encoding & \file{model.py} & 15 lines \\
\textcolor{impactmed}{\textbf{9}} & Normalize target before MSE & \file{loss.py} & 5 lines \\
\textcolor{impactmed}{\textbf{10}} & CLS global token node & \file{model.py} & 30 lines \\
\textcolor{impactmed}{\textbf{11}} & Curriculum masking schedule & \file{dataset.py}, \file{train.py} & 15 lines \\
\textcolor{impactmed}{\textbf{12}} & Multi-resolution windows & \file{dataset.py} & 15 lines \\
\textcolor{impactlow}{\textbf{13}} & Scale up (d\_model=256, layers=4) & \file{config.py} & 1 line \\
\textcolor{impactlow}{\textbf{14}} & 2-layer MLP decoder & \file{model.py} & 5 lines \\
\textcolor{impactlow}{\textbf{15}} & Adaptive K for k-NN graph & \file{model.py} & 3 lines \\
\bottomrule
\end{tabular}
\caption{Prioritized list of improvements. Bugs (red) should be fixed first,
then high-impact changes, then medium.}
\end{table}

\subsection*{Recommended implementation order}

\begin{enumerate}[leftmargin=*]
    \item Fix bugs 1 and 2 immediately --- they affect the correctness of every experiment.
    \item Vectorize \texttt{build\_graph()} --- this is a prerequisite for fast iteration.
    \item Switch to span masking and block dropout augmentation --- these are self-contained changes to \file{dataset.py} with no model changes required.
    \item Add HRV prediction head --- this directly ties pretraining to the clinical features most predictive in downstream tasks.
    \item Add future beat prediction --- requires one new forward pass mode and a small head.
    \item Then consider architectural changes (CLS node, positional encoding, scale-up) once the above objectives are stable.
\end{enumerate}

\medskip
\noindent\textbf{Expected impact of items 1--5 together:}
Better reconstruction targets + harder masking + physiologically meaningful auxiliary objectives
should produce representations that transfer more effectively to HRV-based downstream tasks
(arrhythmia detection, stress, sleep staging) compared to the current BYOL-only global objective.

\begin{thebibliography}{9}
\bibitem{baevski2022data2vec}
Baevski, A., Hsu, W., Xu, Q., Babu, A., Gu, J., \& Auli, M. (2022).
\textit{data2vec: A general framework for self-supervised learning in speech, vision and language}.
ICML 2022.

\bibitem{he2022mae}
He, K., Chen, X., Xie, S., Li, Y., Dollar, P., \& Girshick, R. (2022).
\textit{Masked autoencoders are scalable vision learners}.
CVPR 2022.

\bibitem{grill2020byol}
Grill, J.B., Strub, F., Altch\'e, F., et al. (2020).
\textit{Bootstrap Your Own Latent: A new approach to self-supervised learning}.
NeurIPS 2020.

\bibitem{spanbert}
Joshi, M., Chen, D., Liu, Y., Weld, D., Zettlemoyer, L., \& Levy, O. (2020).
\textit{SpanBERT: Improving pre-training by representing and predicting spans}.
TACL 2020.
\end{thebibliography}

\end{document}
